\subsection{Grothendieck duality over a Macaulay ring}

We continue to assume that A is a local Macaulay ring, of dimension \(d\) . Let \(\Omega\) be a dualizing A-module. The A-module \(\mathrm{H}_{\mathrm{A}}^{d}(\Omega)\) is a Matlis module ( \(\S 8\) , no. 5, Example 6); in accordance with the notation of \(\S 8\) , we set \(\mathrm{D(M)} = \mathrm{Hom}_{\mathrm{A}}(\mathrm{M}, \mathrm{H}_{\mathrm{A}}^{d}(\Omega))\) for every A-module M.

Consider the isomorphism \(\tau^{d}(\Omega):\Omega\otimes\mathrm{H}_{\mathrm{A}}^{d}(\mathrm{A})\to\mathrm{H}_{\mathrm{A}}^{d}(\Omega)\) (no. 2). From this one deduces a homomorphism \(\omega:\mathrm{H}_{\mathrm{A}}^{d}(\mathrm{A})\to\mathrm{Hom}_{\mathrm{A}}(\Omega,\mathrm{H}_{\mathrm{A}}^{d}(\Omega))\) which associates with each element \(u\) of \(\mathrm{H}_{\mathrm{A}}^{d}(\mathrm{A})\) the homomorphism \(x\mapsto\mathrm{H}_{\mathrm{A}}^{d}(f_{x})(u)\) (Remark 3 above).

PROPOSITION 2.--- Let A be a local Macaulay ring of dimension d, and let \(\Omega\) be a dualizing A-module. The homomorphism \(\omega: H_{A}^{d}(A) \to D(\Omega)\) is bijective.

The homomorphism \(\omega\) is the limit of the inductive system of maps \((\omega_{\mathfrak{a}})_{\mathfrak{a}\in\mathcal{D}_{cs}}\) , where

\[
\omega_ {\mathfrak {a}}: \operatorname{Ext} _ {\mathrm{A}} ^ {d} (\mathrm{A} / \mathfrak {a}, \mathrm{A}) \longrightarrow \operatorname{Hom} _ {\mathrm{A}} (\Omega , \operatorname{Ext} _ {\mathrm{A}} ^ {d} (\mathrm{A} / \mathfrak {a}, \Omega))
\]

which associates with each element \(u\) of \(\operatorname{Ext}_{\mathrm{A}}^{d}(\mathrm{A}/\mathfrak{a},\mathrm{A})\) the map \(x\mapsto f_{x}\circ u\) (A, X, p.114). It therefore suffices to prove that each of the maps \(\omega_{\mathfrak{a}}\) is bijective.

Let \(\mathfrak{a}\) be an ideal of \(\mathcal{D}_{cs}\) generated by a sequence \(\mathbf{x} = (x_1, \ldots, x_d)\) completely secant for A. The Koszul complex \(\mathbf{K}^{\bullet}(\mathbf{x}, \mathrm{A})\) provides a projective resolution of A/ \(\mathfrak{a}\) ; for every A-module M, the A-module \(\mathrm{H}^d\big(\mathrm{Homgr}_\mathrm{A}(\mathbf{K}^{\bullet}(\mathbf{x}, \mathrm{A}), \mathrm{M})\big)\) is canonically identified with M/ \(\mathfrak{a}\mathrm{M}\) (A, X, p.155). From this we deduce an isomorphism (A, X, p.100)

\[
\varphi_ {\mathrm{M}}: \mathrm{M} / \mathfrak {a} \mathrm{M} \longrightarrow \operatorname{Ext} ^ {d} (\mathrm{A} / \mathfrak {a}, \mathrm{M}).
\]

Let \(x \in \Omega\) . In view of loc. cit., p.103, prop. 2, we have a commutative diagram

\[
\begin{array}{c c c} \mathrm{A/a} & \xrightarrow {\varphi_ {\mathrm{A}}} & \operatorname{Ext} _ {\mathrm{A}} ^ {d} (\mathrm{A/a}, \mathrm{A}) \\ \bar {f} _ {x} \Bigg \downarrow & & \Bigg \downarrow \operatorname{Ext} (1 _ {\mathrm{A} / a}, f _ {x}) \\ \Omega / a \Omega & \xrightarrow {\varphi_ {\Omega}} & \operatorname{Ext} _ {\mathrm{A}} ^ {d} (\mathrm{A/a}, \Omega) \end{array}
\]

where \(\bar{f}_x\) is the homomorphism deduced from \(f_x\) by passing to quotients. It follows that if for every A-module M one identifies \(\mathrm{Ext}_{\mathrm{A}}^{d}(\mathrm{A}/\mathfrak{a},\mathrm{M})\) with M/ \(\mathfrak{a}\) M by means of \(\varphi_{\mathrm{M}}\) , the homomorphism \(\omega_{\mathfrak{a}}\) is identified with the A-linear map from A/ \(\mathfrak{a}\) to \(\mathrm{Hom}_{\mathrm{A}}(\Omega,\Omega/\mathfrak{a}\Omega)\) which sends 1 to the canonical surjection, that is, to the canonical map A/ \(\mathfrak{a}\longrightarrow\mathrm{End}_{\mathrm{A}/\mathfrak{a}}(\Omega/\mathfrak{a}\Omega)\) . But since the A/ \(\mathfrak{a}\) -module \(\Omega/\mathfrak{a}\Omega\) is dualizing (\(\S\) 9, no. 2, prop. 4), this map is bijective ( \(\S\) 9, no. 4, prop. 6), which proves the proposition.

Identify the Matlis double dual D(D(Ω)) with Ω via the isomorphism \(\widehat{\alpha}_{\Omega}\) (§ 8, no. 3, th. 2, b)).

COROLLARY.--- The homomorphism D(ω): \(\widehat{\Omega} \rightarrow \mathrm{D}(\mathrm{H}_{\mathrm{A}}^{d}(\mathrm{A}))\) is an isomorphism.

Let M be an A-module, and let i be an integer. Consider the canonical homomorphisms (§ 8, no. 7)

\[
\rho_ {d - i} (\mathrm{M}, \Omega): \mathrm{Tor} _ {d - i} ^ {\mathrm{A}} (\mathrm{M}, \mathrm{D} (\Omega)) \longrightarrow \mathrm{D} \big (\mathrm{Ext} _ {\mathrm{A}} ^ {d - i} (\mathrm{M}, \Omega) \big)
\]

\[
\theta^ {d - i} (\mathrm{M}, \mathrm{H} _ {\mathrm{A}} ^ {d} (\mathrm{A})): \operatorname{Ext} _ {\mathrm{A}} ^ {d - i} \left(\mathrm{M}, \mathrm{D} \left(\mathrm{H} _ {\mathrm{A}} ^ {d} (\mathrm{A})\right)\right) \longrightarrow \mathrm{D} \left(\operatorname{Tor} _ {d - i} ^ {\mathrm{A}} \left(\mathrm{M}, \mathrm{H} _ {\mathrm{A}} ^ {d} (\mathrm{A})\right)\right).
\]

Using the isomorphisms \(\omega : \mathrm{H}_{\mathrm{A}}^{d}(\mathrm{A}) \to \mathrm{D}(\Omega)\) , \(\mathrm{D}(\omega) : \widehat{\Omega} \to \mathrm{D}(\mathrm{H}_{\mathrm{A}}^{d}(\mathrm{A}))\) (corollary 1 of Proposition 2) and \(\tau^{i}(\mathrm{M}) : \mathrm{Tor}_{d-i}^{\mathrm{A}}(\mathrm{M}, \mathrm{H}_{\mathrm{A}}^{d}(\mathrm{A})) \longrightarrow \mathrm{H}_{\mathrm{A}}^{i}(\mathrm{M})\) ( \(n^{\circ}\) 2), we deduce canonical homomorphisms of A-modules

\[
\gamma^ {i} (\mathrm{M}): \mathrm{H} _ {\mathrm{A}} ^ {i} (\mathrm{M}) \longrightarrow \mathrm{D} \left(\mathrm{Ext} _ {\mathrm{A}} ^ {d - i} (\mathrm{M}, \Omega)\right)
\]

\[
\delta^ {i} (\mathrm{M}): \operatorname{Ext} _ {\mathrm{A}} ^ {d - i} (\mathrm{M}, \widehat {\Omega}) \longrightarrow \mathrm{D} \left(\mathrm{H} _ {\mathrm{A}} ^ {i} (\mathrm{M})\right).
\]

THEOREM 1 (Grothendieck duality). Let A be a local Macaulay ring of dimension d, and let Ω be a dualizing A-module.

\begin{enumerate}
\def\labelenumi{\alph{enumi})}
\item
  The A-module \(\mathrm{H}_{\mathrm{A}}^{d}(\Omega)\) is a Matlis module; for every A-module P, let D(P) denote the Matlis dual \(\mathrm{Hom}_{\mathrm{A}}(\mathrm{P},\mathrm{H}_{\mathrm{A}}^{d}(\Omega))\) .
\item
  For every \(\mathrm{A}\)-module of finite type M and any integer \(i\) , the canonical homomorphism
\end{enumerate}

\[
\gamma^ {i} (\mathrm{M}): \mathrm{H} _ {\mathrm{A}} ^ {i} (\mathrm{M}) \longrightarrow \mathrm{D} \left(\operatorname{Ext} _ {\mathrm{A}} ^ {d - i} (\mathrm{M}, \Omega)\right)
\]

is an isomorphism of Artinian A-modules.

\begin{enumerate}
\def\labelenumi{\alph{enumi})}
\setcounter{enumi}{2}
\tightlist
\item
  For every A-module M and every integer i, the canonical homomorphism
\end{enumerate}

\[
\delta^ {i} (\mathrm{M}): \operatorname{Ext} _ {\mathrm{A}} ^ {d - i} (\mathrm{M}, \widehat {\Omega}) \longrightarrow \mathrm{D} \left(\mathrm{H} _ {\mathrm{A}} ^ {i} (\mathrm{M})\right)
\]

is an isomorphism of \(\widehat{\mathrm{A}}\) -modules.

This follows from prop. 6 of § 8, no. 7.

COROLLARY.- Let A be a local Macaulay ring, M a nonzero finitely generated A-module of dimension e. The A-module \(\mathrm{H}_{\mathrm{A}}^{e}(\mathrm{M})\) is nonzero.

By virtue of Remark 2 of No. 1, we may assume that the local ring A is complete. In this case A possesses a dualizing module Ω (§ 9, No. 3, cor. 3 of prop. 6); if \(\mathrm{H}_{\mathrm{A}}^{e}(M)\) is zero, then so is its Matlis dual \(\operatorname{Ext}_{\mathrm{A}}^{d-e}(M, \Omega)\), which contradicts Prop. 3, b) of § 9, No. 1.

Remarks.--- 1) When the A-module M is finitely generated, the \(\widehat{\mathrm{A}}\) -module \(\operatorname{Ext}_{\mathrm{A}}^{d-i}(\mathrm{M}, \widehat{\Omega})\) is identified with \(\widehat{\mathrm{A}} \otimes_{\mathrm{A}} \operatorname{Ext}_{\mathrm{A}}^{d-i}(\mathrm{M}, \Omega)\) (A, X, p.108, prop. 7, c)), and \(\delta^{i}(\mathrm{M})\) can also be obtained by composing \(\mathrm{D}(\gamma^{i}(\mathrm{M}))\) with the biduality isomorphism.

\begin{enumerate}
\def\labelenumi{\arabic{enumi})}
\setcounter{enumi}{1}
\tightlist
\item
  Let \(u: \mathrm{M} \to \mathrm{M}'\) be a homomorphism of A-modules. By Remark 1 of No. 2 and that of § 8, No. 7, the following diagrams are commutative:
\end{enumerate}

\[
\begin{array}{c c c} \mathrm{H} _ {\mathrm{A}} ^ {i} (\mathrm{M}) & \xrightarrow {\gamma^ {i} (\mathrm{M})} & \mathrm{D} (\mathrm{Ext} _ {\mathrm{A}} ^ {d - i} (\mathrm{M}, \Omega)) \\ \mathbb {1} _ {\mathrm{A}} ^ {i} (u) \Bigg \downarrow & & \Bigg \downarrow \mathrm{D} (\mathrm{Ext} (u, 1)) \\ \mathrm{H} _ {\mathrm{A}} ^ {i} (\mathrm{M} ^ {\prime}) & \xrightarrow {\gamma^ {i} (\mathrm{M} ^ {\prime})} & \mathrm{D} (\mathrm{Ext} _ {\mathrm{A}} ^ {d - i} (\mathrm{M} ^ {\prime}, \Omega)) \end{array}
\]

\[
\begin{array}{c c c} \operatorname{Ext} _ {\mathrm{A}} ^ {d - i} (\mathrm{M} ^ {\prime}, \widehat {\Omega}) & \xrightarrow {\delta^ {i} (\mathrm{M} ^ {\prime})} & \mathrm{D} (\mathrm{H} _ {\mathrm{A}} ^ {i} (\mathrm{M} ^ {\prime})) \\ \operatorname{Ext} (u, 1) \Bigg \downarrow & & \Bigg \downarrow \mathrm{D} (\mathrm{H} _ {\mathrm{A}} ^ {i} (u)) \\ \operatorname{Ext} _ {\mathrm{A}} ^ {d - i} (\mathrm{M}, \widehat {\Omega}) & \xrightarrow {\delta^ {i} (\mathrm{M})} & \mathrm{D} (\mathrm{H} _ {\mathrm{A}} ^ {i} (\mathrm{M})) \end{array} .
\]

\begin{enumerate}
\def\labelenumi{\arabic{enumi})}
\setcounter{enumi}{2}
\tightlist
\item
  Let
\end{enumerate}

\[
0 \rightarrow \mathrm{M} ^ {\prime} \rightarrow \mathrm{M} \rightarrow \mathrm{M} ^ {\prime \prime} \rightarrow 0\tag{ℰ}
\]

an exact sequence of A-modules. By Remark 2 of No. 2 and that of § 8, No. 7, the following diagrams are commutative:

\[
\begin{array}{c c c} \mathrm{H} _ {\mathrm{A}} ^ {i - 1} (\mathrm{M} ^ {\prime \prime}) & \xrightarrow {\gamma^ {i - 1} (\mathrm{M} ^ {\prime \prime})} & \mathrm{D} (\mathrm{Ext} _ {\mathrm{A}} ^ {d - i + 1} (\mathrm{M} ^ {\prime \prime}, \Omega)) \\ \partial^ {i - 1} (\mathscr {E}) \Bigg \downarrow & & \Bigg \downarrow (- 1) ^ {d - i + 1} \mathrm{D} (\delta^ {d - i} (\mathscr {E}, \Omega)) \\ \mathrm{H} _ {\mathrm{A}} ^ {i} (\mathrm{M} ^ {\prime}) & \xrightarrow {\gamma^ {i} (\mathrm{M} ^ {\prime})} & \mathrm{D} (\mathrm{Ext} _ {\mathrm{A}} ^ {d - i} (\mathrm{M} ^ {\prime}, \Omega)) \end{array}
\]

\[
\begin{array}{c c c} \operatorname{Ext} _ {\mathrm{A}} ^ {d - i} (\mathrm{M} ^ {\prime}, \widehat {\Omega}) & \xrightarrow {\delta^ {i} (\mathrm{M} ^ {\prime})} & \mathrm{D} (\mathrm{H} _ {\mathrm{A}} ^ {i} (\mathrm{M} ^ {\prime})) \\ \delta^ {d - i} (\mathscr {E}, \widehat {\Omega}) \Bigg \downarrow & & \Bigg \downarrow (- 1) ^ {d - i + 1} \mathrm{D} (\partial^ {i - 1} (\mathscr {E})) \\ \operatorname{Ext} _ {\mathrm{A}} ^ {d - i + 1} (\mathrm{M} ^ {\prime \prime}, \widehat {\Omega}) & \xrightarrow {\delta^ {i - 1} (\mathrm{M} ^ {\prime \prime})} & \mathrm{D} (\mathrm{H} _ {\mathrm{A}} ^ {i - 1} (\mathrm{M} ^ {\prime \prime})) \quad . \end{array}
\]

Example.--- Let A be a Noetherian local integral domain of dimension 1; denote by K its field of fractions. Let Ω be a dualizing A-module, and let M be a finitely generated A-module. The A-modules \(\mathrm{H}_{\mathrm{A}}^{0}(M)\) and \(\mathrm{H}_{\mathrm{A}}^{1}(M)\) are canonically identified with T(M) and \((K/A) \otimes_{A} M\) (no. 1, example 5). With these identifications, the duality isomorphisms

\[
\gamma^ {0} (\mathrm{M}): \mathrm{T} (\mathrm{M}) \longrightarrow \mathrm{D} \left(\operatorname{Ext} _ {\mathrm{A}} ^ {1} (\mathrm{M}, \Omega)\right), \quad \gamma^ {1} (\mathrm{M}): (\mathrm{K} / \mathrm{A}) \otimes_ {\mathrm{A}} \mathrm{M} \longrightarrow \mathrm{D} \left(\operatorname{Hom} _ {\mathrm{A}} (\mathrm{M}, \Omega)\right)
\]

(th. 1) are none other than the isomorphisms defined in § 9, no. 6.

